% Source: https://artofproblemsolving.com/community/c3h1547191_geometry
% Stable identity: aops-topic-1547191-q01
Here $M$ is the midpoint of $BC$, $N$ is the midpoint of $AB$, and $O$ is the centroid, which lies two thirds of the way along each median from the vertex. The dashed segment $BO$ is drawn for the argument.

\begin{asy}
usepackage("amsmath");
size(6cm);
defaultpen(fontsize(10pt));
pair A=(1.2,3.4),B=(0,0),C=(5,0);
pair M=(B+C)/2,Nn=(A+B)/2,O=(A+B+C)/3;
fill(A--O--Nn--cycle,paleblue);
draw(A--B--C--cycle,linewidth(1));
draw(A--M,linewidth(0.8));draw(C--Nn,linewidth(0.8));
draw(B--O,dashed+gray(0.4));
dot(M);dot(Nn);dot(O);
label("$A$",A,N);label("$B$",B,SW);label("$C$",C,SE);
label("$M$",M,S);label("$N$",Nn,W);label("$O$",O,dir(50));
label("\mbox{}",(min(currentpicture,user=true)+max(currentpicture,user=true))/2,black); // keeps the text after the figure black when LaTeX typesets the labels inline
\end{asy}

\textbf{Step 1.} Triangles $AOB$ and $ABC$ share the base $AB$. Since $CO=\dfrac{2}{3}\,CN$, the distance from $O$ to line $AB$ is $\dfrac{1}{3}$ of the distance from $C$ to $AB$ (as $ON=\dfrac{1}{3}\,CN$ and $N$ lies on $AB$). So $[AOB]=\dfrac{1}{3}[ABC]$.

\textbf{Step 2.} In triangle $AOB$, $ON$ is a median ($N$ is the midpoint of $AB$), so it splits the triangle into two triangles with equal bases $AN=NB$ and the same height: $[AON]=\dfrac{1}{2}[AOB]$.

Therefore
\[
[AON]=\dfrac{1}{2}\cdot\dfrac{1}{3}[ABC]=\boxed{\dfrac{1}{6}}\,[ABC].
\]
(In fact the three medians cut the triangle into six triangles of equal area, and $AON$ is one of them.)
